% Problem record op_17bba2d8f41697f6. This TeX file is derived from
% database/problems_json/op_17bba2d8f41697f6.json by scripts/sync-tex.mjs; edit the JSON record
% and rerun the script rather than editing this file.
\section{Quantum query complexity of Welded Tree path-finding}
\label{sec:op_17bba2d8f41697f6}

\paragraph{Problem.}
Does finding an explicit ENTRANCE-to-EXIT path in the standard Welded Tree
oracle problem require exponentially many quantum queries?

For a positive integer $n$, construct a graph $G_n$ from two complete binary
trees of depth $n$.  Designate the root of the left tree as ENTRANCE $s$ and
the root of the right tree as EXIT $t$.  Join the $2^n$ leaves of the left
tree and the $2^n$ leaves of the right tree into a uniformly random cycle
that alternates between leaves of the two trees.  Assign the vertices
distinct random binary labels of length $2n$.

The algorithm is given the label of $s$ and coherent quantum-query access to
an adjacency-list oracle.  On a valid vertex label, the oracle returns the
labels of its neighboring vertices, with distinguished invalid outputs used
when necessary; querying arbitrary binary strings therefore does not reveal
global information about the graph unless the corresponding vertex labels
have first been discovered.

The path-finding task is to output an explicit sequence of vertex labels
$v_0,v_1,\ldots,v_\ell$ satisfying
\begin{equation}
  v_0=s,
  \qquad
  v_\ell=t,
  \qquad
  \{v_{j-1},v_j\}\in E(G_n)
  \quad
  \text{for every }1\leq j\leq\ell .
  \label{eq:17bb-weld-path}
\end{equation}
Let $Q_{\mathrm{WTPath}}(n)$ denote the minimum number of oracle queries
used by a quantum algorithm that outputs a valid path satisfying
\eqref{eq:17bb-weld-path} with bounded error on the standard random Welded
Tree oracle distribution.  The open question is whether there exists a
constant $c>0$ such that, for all sufficiently large $n$,
\begin{equation}
  Q_{\mathrm{WTPath}}(n)
  \geq
  2^{cn}.
  \label{eq:17bb-weld-exponential}
\end{equation}
Equivalently, does the exponential quantum-query lower bound in
\eqref{eq:17bb-weld-exponential} hold for unrestricted quantum algorithms?

\subsection*{Status}

Unsolved

\subsection*{Source}

Explicitly posed by Scott Aaronson in Aaronson (2021), Problem 10, as the
question of whether a quantum computer requires exponentially many queries
to find a left-to-right path in the Welded Tree oracle problem
\sourcecite{ref:17bb-aaronson}{Aar21}.  The distinction originates in the
Welded Tree algorithm of Childs et al. (2003), whose quantum walk finds the
EXIT but does not identify a particular ENTRANCE-to-EXIT path; Childs,
Coudron, and Gilani (2023) describe efficient path-finding versus quantum
hardness as a longstanding open problem and prove a restricted no-go result.
The statement above is a self-contained formulation of Aaronson's question
rather than a verbatim quotation.

\subsection*{Progress}

\begin{itemize}
  \item Childs, Cleve, Deotto, Farhi, Gutmann, and Spielman introduced the Welded
Tree oracle problem as an early example of an exponential quantum speedup
based on quantum walks rather than the quantum Fourier transform: their
continuous-time quantum walk reaches the EXIT using polynomially many oracle
queries, whereas classical randomized algorithms require exponentially many
queries in the relevant black-box model.  Crucially, the quantum algorithm
finds the label of the EXIT without producing a particular path from the
ENTRANCE to the EXIT
\sourcecite{ref:17bb-childs}{CCGS03}.

  \item The graph always contains short ENTRANCE-to-EXIT paths, so the difficulty
is not an exponentially long required output; rather, the quantum walk
propagates amplitude through the graph by interference without retaining
enough classical information to reconstruct which sequence of vertices
constitutes such a path
\sourcecite{ref:17bb-childs}{CCGS03}.

  \item Aaronson isolated the question explicitly as Problem 10 in his 2021 list
of open problems in quantum query complexity, asking whether even a quantum
computer requires exponentially many queries to find a left-to-right path in
the Welded Tree graph; the lower-bound formulation, rather than merely the
absence of a known polynomial-time algorithm, is the canonical form of the
open problem
\sourcecite{ref:17bb-aaronson}{Aar21}.

  \item Childs, Coudron, and Gilani introduced a class of genuine rooted quantum
algorithms: within each branch of their superposition they always store a
set of vertex labels forming a connected subgraph that includes the
ENTRANCE, and they only provide such meaningful vertex labels as oracle
inputs.  They proved that these genuine rooted quantum algorithms cannot
find an ENTRANCE-to-EXIT path with polynomially many queries
\sourcecite{ref:17bb-ccg}{CCG23}.

  \item The restricted result does not establish the exponential lower bound for
unrestricted algorithms: the EXIT-finding quantum walk is itself not rooted
in the relevant sense, since it can discard information about the route by
which its amplitude reached a vertex.  A general path-finding algorithm
might therefore exploit precisely this non-rooted, forgetful behavior, and
hardness proofs require techniques that can handle such algorithms
\sourcecite{ref:17bb-ccg}{CCG23}.

  \item Efficient quantum path-finding has subsequently been demonstrated on
graph families inspired by Welded Trees: Li and Zur constructed a regular
welded-tree circuit graph for which multidimensional quantum walks and
electrical-flow states yield an exponential quantum speedup for finding an
explicit path.  This result concerns a modified graph family and does not
give an efficient path-finding algorithm for the original randomly welded
binary-tree oracle
\sourcecite{ref:17bb-li-zur}{LZ25}.
\end{itemize}

\subsection*{References}

\begin{enumerate}[label={},leftmargin=4.8em,labelsep=0.5em]
  \footnotesize
  \item[\textup{[CCGS03]}]\label{ref:17bb-childs}
  Andrew M. Childs, Richard Cleve, Enrico Deotto, Edward Farhi, Sam
Gutmann, and Daniel A. Spielman, ``Exponential Algorithmic Speedup by Quantum
Walk,'' in \emph{Proceedings of the 35th Annual ACM Symposium on Theory of
Computing (STOC 2003)}, 59--68 (2003).
\href{https://doi.org/10.1145/780542.780552}{doi:10.1145/780542.780552};
\href{https://arxiv.org/abs/quant-ph/0209131}{arXiv:quant-ph/0209131}.

  \item[\textup{[Aar21]}]\label{ref:17bb-aaronson}
  Scott Aaronson, ``Open Problems Related to Quantum Query Complexity,'' \emph{ACM
Transactions on Quantum Computing} \textbf{2}, Article 14 (2021).
\href{https://doi.org/10.1145/3488559}{doi:10.1145/3488559};
\href{https://arxiv.org/abs/2109.06917}{arXiv:2109.06917}.

  \item[\textup{[CCG23]}]\label{ref:17bb-ccg}
  Andrew M. Childs, Matthew Coudron, and Amin Shiraz Gilani, ``Quantum
Algorithms and the Power of Forgetting,'' in \emph{14th Innovations in
Theoretical Computer Science Conference (ITCS 2023)}, LIPIcs \textbf{251},
37:1--37:22 (2023).
\href{https://doi.org/10.4230/LIPIcs.ITCS.2023.37}{doi:10.4230/LIPIcs.ITCS.2023.37};
\href{https://arxiv.org/abs/2211.12447}{arXiv:2211.12447}.

  \item[\textup{[LZ25]}]\label{ref:17bb-li-zur}
  Jianqiang Li and Sebastian Zur, ``Multidimensional Electrical Networks
and their Application to Exponential Speedups for Graph Problems,'' \emph{Quantum}
\textbf{9}, 1733 (2025).
\href{https://doi.org/10.22331/q-2025-05-06-1733}{doi:10.22331/q-2025-05-06-1733};
\href{https://arxiv.org/abs/2311.07372}{arXiv:2311.07372}.
\end{enumerate}

\subsection*{Comment}

Two qualitatively different alternatives remain for the original Welded
Tree oracle: either an unrestricted exponential quantum-query lower bound
of the form \eqref{eq:17bb-weld-exponential}, extending beyond the
genuine-rooted model, or a quantum algorithm that finds an explicit path
using subexponentially many queries.
A polynomial-query path-finding algorithm would show an even stronger
failure of the exponential conjecture in
\eqref{eq:17bb-weld-exponential}.

\subsection*{Field}

Quantum algorithm

\subsection*{Topic}

Computational complexity and computability

\subsection*{ID}

\texttt{op\_17bba2d8f41697f6}
